\documentclass[a4paper, 11pt]{article}

% --- Packages ---
\usepackage[utf8]{inputenc}
\usepackage[T1]{fontenc}
% \usepackage[french]{babel}
\usepackage[margin=1.5cm, top=2cm, bottom=2cm]{geometry}
\usepackage{xcolor}
\usepackage{fancyhdr}
\usepackage{titlesec}
\usepackage{enumitem}
\usepackage{array}
\usepackage{longtable}
\usepackage{booktabs}
\usepackage{colortbl}

% --- Définition des couleurs ---
\definecolor{primaryBlue}{RGB}{20, 50, 90}
\definecolor{lightBlue}{RGB}{235, 240, 247}

% --- Styles de document ---
\pagestyle{fancy}
\fancyhf{}
\fancyfoot[L]{\footnotesize \textcolor{gray}{Projet Wakala - Feuille de Route}}
\fancyfoot[R]{\footnotesize \textcolor{gray}{Page \thepage}}
\renewcommand{\headrulewidth}{0pt}
\renewcommand{\footrulewidth}{0.4pt}

% --- Personnalisation des titres ---
\titleformat{\section}{\Large\bfseries\color{primaryBlue}}{}{0em}{}[\titlerule]
\titleformat{\subsection}{\large\bfseries\color{primaryBlue}}{}{0em}{}

\newcolumntype{L}[1]{>{\raggedright\arraybackslash}p{#1}}

\begin{document}

% --- En-tête ---
\begin{center}
    {\Huge \textcolor{primaryBlue}{\textbf{WAKALA}}}\\[0.2cm]
    {\large \bfseries \textcolor{primaryBlue}{Feuille de Route du Stage}}
\end{center}

\vspace{0.5cm}

% --- Introduction ---
\section*{Introduction}
Ce document présente le calendrier et la planification (Feuille de Route) du projet de stage. L'objectif global est le développement de bout en bout d'une plateforme d'accompagnement à l'achat automobile, structurée autour de deux axes fondamentaux : un Chatbot interactif et un Moteur de Recommandation IA. Ce planning permet de garantir l'avancée structurée du projet.

\vspace{0.3cm}

% --- Contenu ---
\section*{Contenu}

\subsection*{Phase 1 : Cadrage et Développement du Chatbot}
Durant cette première étape, la priorité est donnée à la création de l'interface conversationnelle et à la compréhension des requêtes de l'utilisateur.
\begin{itemize}[label=\textbullet]
    \item \textbf{Étude et Architecture :} Définition des outils de base pour développer l'application.
    \item \textbf{Intégration de l'Intelligence Artificielle :} Mise en place du cerveau du Chatbot pour assurer des réponses de qualité en toute confidentialité. Le Chatbot devra être capable de poser des questions de clarification (ex: "Avez-vous besoin d'un grand coffre ?").
    \item \textbf{Compréhension des Intentions :} Développement du système permettant au Chatbot de comprendre le budget, les envies (sportivité, économie) et les craintes depuis une phrase naturelle écrite par l'utilisateur, pour générer un profil acheteur précis.
\end{itemize}

\vspace{0.2cm}
\subsection*{Phase 2 : Développement du Moteur de Recommandation}
Cette phase est purement algorithmique. Il s'agit de construire le moteur capable de lier le profil de l'utilisateur aux véhicules.
\begin{itemize}[label=\textbullet]
    \item \textbf{Préparation des données :} Rassemblement et nettoyage des annonces de véhicules.
    \item \textbf{Création des Profils :} Traduction mathématique des caractéristiques des voitures (consommation, fiabilité, espace) et des envies de l'acheteur.
    \item \textbf{Base de données Intelligente :} Configuration du système de stockage permettant des recherches instantanées parmi des milliers d'annonces.
    \item \textbf{Logique de Recommandation :} Développement de la formule pour trouver les meilleures voitures pour un profil donné. Le moteur apprendra que si un acheteur demande une "voiture économique pour la ville", il faut privilégier les citadines essence ou hybrides avec un faible kilométrage, même si l'acheteur n'a pas précisé de modèle.
\end{itemize}

\vspace{0.2cm}
\newpage
\subsection*{Phase 3 : Intégration et Tests}
Cette dernière étape vise à assembler le système et à valider l'expérience utilisateur finale.
\begin{itemize}[label=\textbullet]
    \item \textbf{Fusion des Systèmes :} Connexion du Chatbot au Moteur de Recommandation pour une expérience fluide.
    \item \textbf{Ajustement Fin :} Optimisation des suggestions pour correspondre exactement au profil de l'utilisateur.
    \item \textbf{Tests Utilisateurs (UX) :} Vérification que le client se sent bien accompagné et guidé.
\end{itemize}

\vspace{0.2cm}
\subsection*{Phase 4 : Déploiement et Évaluation Continue}
Cette ultime étape consiste à mettre la plateforme en production et à monitorer ses performances réelles auprès des utilisateurs cibles.
\begin{itemize}[label=\textbullet]
    \item \textbf{Mise en Production :} Déploiement de l'application sur une infrastructure cloud performante pour garantir sa disponibilité.
    \item \textbf{Suivi des Métriques (Analytics) :} Mise en place de tableaux de bord pour analyser le taux de conversion et l'engagement des utilisateurs.
    \item \textbf{Amélioration Continue :} Collecte des retours pour corriger d'éventuels bugs, enrichir les données du moteur et affiner les réponses de l'IA.
\end{itemize}

\vspace{0.3cm}

% --- Synthèse ---
\section*{Synthèse}
Ce projet s'articule autour de quatre grandes phases, conçues pour assurer une progression claire et logique. Il débute par la compréhension des besoins de l'utilisateur (Chatbot), se poursuit par l'analyse du marché (Moteur de Recommandation), puis passe par l'intégration technique de ces briques. Enfin, la dernière phase est consacrée au déploiement et à l'optimisation de l'outil. Cette démarche étape par étape permet de respecter le planning tout en assurant la réussite globale de ce projet d'Intelligence Artificielle.

\end{document}
